\subsection{Fibadic 环境塌缩、appearance-rank 原始谱与循环受体阈值}\label{subsec:conclusion-fibadic-packet-primitive-spectrum-cyclic-threshold}
沿用小节 \ref{subsec:conclusion-fibadic-profinite-collapse-character-conductor-rigidity}、
\ref{subsec:conclusion-fibadic-new-harmonics-conductor-cyclotomic}、
\ref{subsec:conclusion-fibadic-gcd-operator-bernoulli-ramanujan}、
\ref{subsec:conclusion-mixed-hiddenstate-finite-abelian-observability}
与
\ref{subsec:conclusion-fibadic-kernel-semisimplicity-hiddenstate-truncation}
中的记号。此前结论已经分别建立了 Fibonacci-adic 地址完成向
\(\widehat{\ZZ}\) 的全局塌缩、有限值连续观测的唯一 ordinary conductor、
精确深度新空间 \(\mathcal N_d\) 的 Möbius 投影结构、cyclotomic 包
\(\Pi_d\) 与 Ramanujan 包核 \(K_d\) 的对应，以及混合隐藏态
\[
\mathcal H_S\cong (\ZZ/2\ZZ)^\beta\times \widehat X_F
\]
的有限阿贝尔观测分类。把这些链条进一步并置之后，可将 fibadic 侧的全部专属性压缩为一条更紧的终端结构：ambient profinite 几何与连续 \(1\)-保持环对称完全普通化，而真正的新信息则严格落在 exact-depth packet、其 primitive spectrum 与非循环有限受体阈值上。

\begin{theorem}[环境塌缩与精确深度模态饱和的并存]\label{thm:conclusion-fibadic-collapse-packet-saturation-coexistence}
成立规范同构
\[
\widehat X_F\cong \widehat{\ZZ},
\qquad
\operatorname{Aut}_{\mathrm{cont},1\text{-ring}}(\widehat X_F)=\{\mathrm{id}\},
\]
但与此同时，精确深度新模态数与 cyclotomic 包次数仍满足
\[
a_d=F_d+O\!\bigl(\tau(d)\varphi^{d/2}\bigr),
\qquad
\deg \Pi_d=F_d+O\!\bigl(\tau(d)\varphi^{d/2}\bigr),
\]
从而
\[
\frac{a_d}{F_d}\longrightarrow 1,
\qquad
\frac{\deg \Pi_d}{F_d}\longrightarrow 1.
\]
因此，fibadic 体系的专属性并不残留于 ambient profinite 对象或连续 \(1\)-保持环自同构，而几乎全部转移到 exact-depth 过滤与其对应的新包分解之中。
\end{theorem}
\begin{proof}
拓扑环塌缩由定理
\ref{thm:conclusion-fibadic-profinite-collapse-to-zhat}
给出，而连续 \(1\)-保持环自同构刚性由定理
\ref{thm:conclusion-fibadic-global-continuous-endomorphism-classification}
给出。另一方面，\(a_d\) 的渐近公式与比值极限见推论
\ref{cor:conclusion-fibadic-exact-depth-characters-dominate}；
\(\deg\Pi_d\) 的相应结论见推论
\ref{cor:conclusion-fibadic-new-harmonics-dominate}。
合并即得全部断言。证毕。
\end{proof}

\begin{theorem}[有限值连续观测的导体--深度传递]\label{thm:conclusion-fibadic-observable-conductor-depth-transfer}
设 \(Y\) 为有限离散集，\(v:\widehat X_F\to Y\) 为连续映射。记
\[
N_v
\]
为定理
\ref{thm:conclusion-fibadic-finite-valued-observable-conductor}
给出的唯一 ordinary conductor，并定义其最小 Fibonacci 分辨率指数
\[
R(v):=
\min\bigl\{
r\ge 0:\ \exists\,\widetilde v_r:X_r\to Y
\ \text{使}\ v=\widetilde v_r\circ \rho_r
\bigr\},
\]
其中 \(\rho_r:\widehat X_F\twoheadrightarrow X_r\) 为标准截断映射。则
\[
R(v)=d_{N_v}-2,
\qquad
d_{N_v}
=
\operatorname{lcm}_{p^e\parallel N_v} z(p^e).
\]
特别地，任意两个具有同一 conductor 的有限值连续观测，必在同一最小 Fibonacci 分辨率上首次整体可见。
\end{theorem}
\begin{proof}
由定理 \ref{thm:conclusion-fibadic-finite-valued-observable-conductor}，
\(v\) 唯一经由
\[
\pi_{N_v}:\widehat X_F\twoheadrightarrow \ZZ/(N_v\ZZ)
\]
因子化。又由定理
\ref{thm:conclusion-fibadic-minimal-resolution-for-prime-register-package}，
模 \(N_v\) 读出经由 \(X_r\) 因子化，当且仅当
\[
d_{N_v}\mid (r+2).
\]
故最小此类 \(r\) 恰为 \(d_{N_v}-2\)。最后一句只是这一公式对同 conductor 观测的直接推论。证毕。
\end{proof}

\begin{theorem}[可见深度条件期望代数的原始谱恰为 appearance ranks]\label{thm:conclusion-fibadic-visible-algebra-primitive-spectrum}
\leanverified{paper\_conclusion\_fibadic\_visible\_algebra\_primitive\_spectrum}
记
\[
\mathcal A_{\mathrm{vis}}
:=
\operatorname{span}\{E_d:\ d\ge 1\},
\]
并称满足 \(a_n>0\) 的整数 \(n\) 为一个 appearance rank。则 \(\mathcal A_{\mathrm{vis}}\) 是交换半单代数，并具有规范分解
\[
\mathcal A_{\mathrm{vis}}
\cong
\bigoplus_{n:\,a_n>0}\CC\,P_n,
\qquad
T_\alpha=\sum_{n\ge 1}\widehat\alpha(n)P_n,
\]
其中 \(P_n\) 为精确深度新空间 \(\mathcal N_n\) 上的原始中心幂等元，且
\[
\widehat\alpha(n)=\sum_{n\mid d}\alpha(d).
\]
从而：
\begin{enumerate}
  \item primitive spectrum 满足
  \[
  \operatorname{Prim}(\mathcal A_{\mathrm{vis}})
  \cong
  \{n\ge 1:\ a_n>0\};
  \]
  \item exact depth 不只是一个分级参数，而是 \(\mathcal A_{\mathrm{vis}}\) 的原始谱点；
  \item 任意把两个不同 appearance ranks \(m\neq n\) 识别为同一点的交换商，都会在某个有限深度可见算子上失去忠实性。
\end{enumerate}
\end{theorem}
\begin{proof}
定理
\ref{thm:conclusion-fibadic-visible-depth-mobius-primitive-semisimplicity}
已给出
\[
\mathcal A_{\mathrm{vis}}
\cong
c_{00}\bigl(\{n\ge 1:\ a_n>0\}\bigr)
\]
的点乘代数模型，并把 \(P_n\) 识别为坐标原始幂等元。故 primitive ideals 恰为各坐标评价
\[
\operatorname{ev}_n:\mathcal A_{\mathrm{vis}}\to \CC,
\qquad
\operatorname{ev}_n(T_\alpha)=\widehat\alpha(n),
\]
的核，这就证明了前两条。

若某交换商把 \(m\neq n\) 识别为同一点，则该商中的任意元素在这两点上的坐标值必须相同。取一个有限支撑标量族 \(\lambda\) 满足
\[
\lambda(m)=1,
\qquad
\lambda(n)=0.
\]
由定理
\ref{thm:conclusion-fibadic-visible-depth-mobius-primitive-semisimplicity}
中的 Möbius 反演，存在唯一 \(T_\alpha\in \mathcal A_{\mathrm{vis}}\) 使
\[
T_\alpha|_{\mathcal N_k}=\lambda(k)\,\mathrm{id}_{\mathcal N_k}.
\]
该算子在 \(m\) 与 \(n\) 上的作用不同，故其像在上述商中不可能仍然忠实。证毕。
\end{proof}

\begin{theorem}[精确深度包的 cyclotomic、Bernoulli 与 Ramanujan 三重完备性]\label{thm:conclusion-fibadic-depth-packet-triple-completeness}
对每个 \(d>1\)，定义
\[
R_d:=\{N\ge 1:\ d_N=d\},
\]
\[
\Pi_d(X)=\prod_{N\in R_d}\Phi_N(X),
\qquad
J_d(u):=\log \Pi_d(e^u),
\qquad
K_d(x):=\sum_{N\in R_d}c_N(\pi_N(x)).
\]
则上述三类对象彼此完全等价，并且每一类都唯一决定其余两类。更精确地说：
\begin{enumerate}
  \item \(\Pi_d\) 的 cyclotomic 分解唯一恢复 \(R_d\)；
  \item \(J_d\) 作为 \(u=0\) 邻域中的解析 germ 唯一恢复 \(\Pi_d\)，故也唯一恢复 \(R_d\)；
  \item \(K_d\) 的 Fourier 支撑恰为全部满足 \(\delta(\chi)=d\) 的连续特征，因此它唯一恢复 \(R_d\)。
\end{enumerate}
\end{theorem}
\begin{proof}
由定理 \ref{thm:conclusion-fibadic-cyclotomic-packet-decomposition}，
\[
\Pi_d(X)=\prod_{N\in R_d}\Phi_N(X),
\]
故 \(\Pi_d\) 的 cyclotomic 因子分解唯一确定 \(R_d\)。

对 \(d>1\)，\(\Pi_d(1)\in \NN_{>0}\)，故 \(J_d(u)=\log \Pi_d(e^u)\) 在 \(u=0\) 邻域有唯一实值解析支。由 \(J_d\) 可恢复
\[
\Pi_d(e^u)=\exp(J_d(u))
\]
在 \(u=0\) 邻域的解析 germ；再由 \(X=e^u\) 的局部可逆性，可恢复 \(\Pi_d(X)\) 在 \(X=1\) 邻域的 germ，而多项式由其任一非空开邻域 germ 唯一决定，故 \(J_d\) 唯一恢复 \(\Pi_d\)，从而恢复 \(R_d\)。

另一方面，定理 \ref{thm:conclusion-fibadic-exact-depth-kernel-ramanujan-packet} 给出
\[
K_d(x)=\sum_{\delta(\chi)=d}\chi(x)
=
\sum_{N\in R_d}c_N(\pi_N(x)).
\]
角色系在 \(L^2(\widehat X_F,m_H)\) 中构成正交基，故 \(K_d\) 唯一决定其 Fourier 支撑
\[
\{\chi:\ \delta(\chi)=d\}.
\]
该支撑中出现的全部普通导体恰为 \(R_d\)。于是 \(K_d\) 也唯一决定 \(R_d\)。既然三者都唯一决定同一集合 \(R_d\)，它们便彼此完全等价。证毕。
\end{proof}

\begin{theorem}[深度包在 \texorpdfstring{$X=1$}{X=1} 处的一点射流证书]\label{thm:conclusion-fibadic-depth-packet-onepoint-jet-certificate}
对每个 \(d>1\)，成立
\[
\Pi_d(1)>1
\iff
\exists\,p^k\ \text{使}\ z(p^k)=d,
\]
\[
\frac{2\Pi_d'(1)}{\Pi_d(1)}
=
K_d(0)
=
\|K_d\|_{L^2(m_H)}^2
=
a_d,
\]
以及
\[
\left.\frac{\dd^2}{\dd u^2}\log \Pi_d(e^u)\right|_{u=0}
=
\frac{S_2(d)}{12}.
\]
进一步，
\[
\frac{2\Pi_d'(1)}{\Pi_d(1)F_d}
=
\frac{a_d}{F_d}
\longrightarrow 1.
\]
因此，深度包在 \(X=1\) 处的一阶射流已经同时锁定“是否出现真正新的素幂频率”与“该层全部可见谱质量”，而二阶射流则恢复 Möbius 加权平方和 \(S_2(d)\)。
\end{theorem}
\begin{proof}
第一条由命题
\ref{prop:conclusion-fibadic-depth-content-prime-power-detection}
与推论
\ref{cor:conclusion-fibadic-cyclotomic-packet-detects-new-prime-power}
立即得到。又由推论
\ref{cor:conclusion-fibadic-bernoulli-jet-kd-zero-lock}，
\[
\frac{2\Pi_d'(1)}{\Pi_d(1)}=K_d(0).
\]
而定理 \ref{thm:conclusion-fibadic-exact-depth-kernel-ramanujan-packet}
给出
\[
K_d(0)=a_d,
\qquad
\|K_d\|_{L^2(m_H)}^2=a_d.
\]
于是得到中间一串恒等式。

二阶射流公式则正是定理
\ref{thm:conclusion-fibadic-depth-packet-bernoulli-mobius-jet}
的特例。最后一式由
\[
\frac{2\Pi_d'(1)}{\Pi_d(1)}=a_d
\]
与推论 \ref{cor:conclusion-fibadic-exact-depth-characters-dominate}
中的
\[
\frac{a_d}{F_d}\to 1
\]
直接推出。证毕。
\end{proof}

\begin{theorem}[有限阿贝尔受体上的双资源正交与预算不可互偿]\label{thm:conclusion-fibadic-abelian-receiver-budget-orthogonality}
记
\[
\mathcal H_S:=K_S^{\mathrm{col}}\times \widehat X_F,
\qquad
\mathcal H_{S,N}:=(\ZZ/2\ZZ)^\beta\times \ZZ/(N\ZZ).
\]
则对任意有限阿贝尔群 \(G\)，都有自然同构
\[
\operatorname{Hom}_{\mathrm{cont}}(\mathcal H_S,G)
\cong
G[2]^{\,\beta}\oplus G.
\]
并且若要求存在忠实嵌入
\[
\mathcal H_{S,N}\hookrightarrow G,
\]
则必有
\[
|G|\ge 2^\beta N,
\]
且等号成立当且仅当
\[
G\cong (\ZZ/2\ZZ)^\beta\times \ZZ/(N\ZZ).
\]
因此，碰撞扭量预算与 ordinary conductor 预算在有限阿贝尔层面严格正交，不存在以额外 \(2\)-扭量资源抵消 conductor 预算的机制。
\end{theorem}
\begin{proof}
第一式正是定理
\ref{thm:conclusion-mixed-hidden-state-finite-abelian-classification}
的内容。最小基数界与极小同构类型则由定理
\ref{thm:conclusion-mixed-hidden-state-minimal-finite-receiver}
直接给出。由这两个结论可见，碰撞方向只经由 \(G[2]\) 进入，而地址方向只经由一个 ordinary finite Abelian 坐标进入，二者在同态分类中形成直和分裂，因而不可互偿。证毕。
\end{proof}

\begin{corollary}[单循环受体只存在于薄例外扇区]\label{cor:conclusion-fibadic-cyclic-receiver-thin-exceptional-sector}
存在忠实嵌入
\[
\mathcal H_{S,N}\hookrightarrow \ZZ/(M\ZZ)
\]
当且仅当满足以下二者之一：
\[
\beta=0,
\qquad\text{或}\qquad
\beta=1\ \text{且}\ N\ \text{为奇数}.
\]
在这两种情形下，最小可能的 \(M\) 分别为
\[
M=N
\quad (\beta=0),
\qquad
M=2N
\quad (\beta=1,\ N\ \text{奇}).
\]
因此，只要
\[
\beta\ge 2,
\qquad\text{或}\qquad
\beta=1\ \text{且}\ 2\mid N,
\]
任何单循环受体都不可能忠实承载混合隐藏态。
\end{corollary}
\begin{proof}
这正是定理
\ref{thm:conclusion-mixed-hiddenstate-cyclic-receiver-criterion}
的直接重述。证毕。
\end{proof}

\begin{conjecture}[fibadic 忠实 \texorpdfstring{$\pi$}{pi}-通道的三寄存器原理]\label{conj:conclusion-fibadic-pi-channel-three-register-principle}
若一个 \(\pi\)-通道要在 fibadic 语境中同时忠实承载地址、exact-depth 新生包与碰撞扭量，则最小外置状态不应退化为单一标量或单一循环寄存器，而至少应同时包含以下三类数据：
\[
N
\quad\text{(ordinary conductor)},
\]
\[
\mathcal P_d
\simeq
\Pi_d
\simeq
J_d
\simeq
K_d
\quad\text{(exact-depth packet)},
\]
\[
\xi\in (\ZZ/2\ZZ)^\beta
\quad\text{(collision torsion)}.
\]
其中，缺失第一项将丧失有限寻址；缺失第二项只能看到旧层代数而看不到当层新包；缺失第三项则对碰撞型 \(\mathsf{NULL}\) 结构性失明。特别地，由推论
\ref{cor:conclusion-fibadic-cyclic-receiver-thin-exceptional-sector}
可知，企图用单循环状态统一忠实承载这三类资源，只可能在薄例外扇区
\[
\beta=0
\qquad\text{或}\qquad
(\beta=1,\ N\ \text{奇})
\]
中成立；在 generic 情形下，非循环外置是不可避免的。
\end{conjecture}

\paragraph{统一读法}
定理 \ref{thm:conclusion-fibadic-collapse-packet-saturation-coexistence} 至推论
\ref{cor:conclusion-fibadic-cyclic-receiver-thin-exceptional-sector}
把 fibadic 侧原本分散的若干链条收束为同一组终端后果：ambient 对象层塌缩为
\(\widehat{\ZZ}\)，连续 \(1\)-保持环对称完全刚化；有限值观测的全部 Fibonacci 专属性压缩为导体到最小可见深度的传递
\[
N\longmapsto d_N;
\]
appearance ranks 不再只是计数标签，而成为可见深度条件期望代数的 primitive spectrum；exact-depth packet 又同时具有 cyclotomic、多项式 jet 与 Ramanujan kernel 三种彼此完备的实现；最后，一旦进入有限阿贝尔受体层，collision torsion 与 conductor 预算便严格正交，而单循环忠实承载只剩下一条薄例外扇区。于是，fibadic 体系真正的新颖性已不在 ambient profinite 几何，而在 exact-depth packet 与 collision torsion 所构成的双层隐藏结构；一旦要求忠实观测，单标量化便几乎处处失败。

\endinput
